% The full cellular diagram: long exact sequences of adjacent skeletons
% strung together, generating the boundary maps d_k (blue).
% Source: book/tex/homology/cellular.tex (§ Cellular chain complex).
\documentclass[border=2mm]{standalone}
\usepackage{amssymb}
\usepackage{amsmath}
\usepackage{tikz-cd}
\newcommand{\CX}[1]{\boxed{\color{blue}\operatorname{Cells}_{#1}(X)}}
\begin{document}
\begin{tikzcd}[column sep=tiny]
  & \underbrace{H_3(X^2)}_{=0} \arrow[d, "0"] \\
  \CX{4} \arrow[r, "\partial_4"] \arrow[rd, "d_4", blue]
    & H_3(X^3) \arrow[r, two heads] \arrow[d, hook]
    & \underbrace{H_3(X^4)}_{\cong H_3(X)} \arrow[r, "0"]
    & \underbrace{H_3(X^4, X^3)}_{= 0} \\
  & \CX{3} \arrow[d, "\partial_3"] \arrow[rd, "d_3", blue]
    && \underbrace{H_1(X^0)}_{=0} \arrow[d, "0"] \\
  \underbrace{H_2(X^1)}_{=0} \arrow[r, "0"]
    & H_2(X^2) \arrow[r, hook] \arrow[d, two heads]
    & \CX{2} \arrow[r, "\partial_2"] \arrow[rd, "d_2", blue]
    & H_1(X^1) \arrow[r, two heads] \arrow[d, hook]
    & \underbrace{H_1(X^2)}_{\cong H_1(X)} \arrow[r, "0"]
    & \underbrace{H_1(X^2, X^1)}_{=0} \\
  & \underbrace{H_2(X^3)}_{\cong H_2(X)} \arrow[d, "0"]
    && \CX{1} \arrow[d, "\partial_1"] \arrow[rd, "d_1", blue] \\
  & \underbrace{H_2(X^3, X^2)}_{=0}
    & \underbrace{H_0(\varnothing)}_{=0} \arrow[r, "0"]
    & H_0(X^0) \arrow[r, hook] \arrow[d, two heads]
    & \CX{0} \arrow[r, "\partial_0"]
    & \dots \\
  &&& \underbrace{H_0(X^1)}_{\cong H_0(X)} \arrow[d, "0"] \\
  &&& \underbrace{H_0(X^1, X^0)}_{=0}
\end{tikzcd}
\end{document}
